%============================================================
% INTRODUCTION GÉNÉRALE
%============================================================
\chapter*{Introduction générale}
\addcontentsline{toc}{chapter}{Introduction générale}
\markboth{Introduction générale}{Introduction générale}


La télématique embarquée permet la collecte, le traitement et la transmission à distance de données issues de véhicules connectés. Face à la diversité des cas d'usage — du simple traceur GPS au système de supervision avancée pour poids lourds — les fabricants sont confrontés à un défi structurel : maintenir des bases de code distinctes pour chaque variante de produit engendre une dette technique croissante et des cycles de qualification allongés. C'est dans ce contexte que s'inscrit le projet \textit{GlobalFleet}, réalisé au sein de \textbf{Smart Automation Technologies} : concevoir et implémenter un firmware embarqué générique, modulaire et configurable sur microcontrôleur \textbf{ESP32-P4} sous \textbf{FreeRTOS}, capable de couvrir trois variantes de produits à partir d'un seul code source, tout en respectant les contraintes de fiabilité, de ressources et de sécurité inhérentes au déploiement industriel sur véhicules.